\section{用例}
\subsection{フォント}

\begin{itemize}
\item 普通の文字(normal)
\item \emph{強調(emph)}
\item \textbf{強い強調(textbf)}
\item \sout{打ち消し線}
\item \textsuperscript{上付き}
\end{itemize}

\subsection{quote環境}

\begin{quote}
吾輩は猫である。名前はまだ無い。

どこで生れたかとんと見当がつかぬ。
\end{quote}

\subsection{verbatim環境}

半角80文字がちょうど収まるサイズにしています。
コマンドラインやソースコードを示す場合はverbatim環境よりもlistingsを推奨します

\begin{verbatim}
01234567890123456789012345678901234567890123456789012345678901234567890123456789
０１２３４５６７８９０１２３４５６７８９０１２３４５６７８９０１２３４５６７８９
\end{verbatim}

フォントが小さくなって嫌だ、という方のために60文字がちょうど収まるverbatim60とverbatim40を用意しています。

\subsection{verbatim60}

\begin{verbatim60}
012345678901234567890123456789012345678901234567890123456789
０１２３４５６７８９０１２３４５６７８９０１２３４５６７８９
\end{verbatim60}

\subsection{verbatim40}

\begin{verbatim40}
0123456789012345678901234567890123456789
０１２３４５６７８９０１２３４５６７８９
\end{verbatim40}

\section{コードリスティング}

\subsection{コマンドライン}
\begin{lstlisting}[language=command]
# uname --machine
x86_64
\end{lstlisting}

\subsection{C}
\begin{lstlisting}[language=C]
#include <stdio.h>

int main(int argc, char *argv[]){
    printf("Hello World!");
    return 0;
}
\end{lstlisting}

\subsection{Python}
\begin{lstlisting}[language=python]
def hello():
    print('Hello World!')

if __name__ == '__main__':
    hello()
\end{lstlisting}

\subsection{80文字テスト}
\begin{lstlisting}
01234567890123456789012345678901234567890123456789012345678901234567890123456789
\end{lstlisting}


\subsection{リスト}

\subsection{列挙リスト}
\begin{itemize}
\item りんご
\item みかん
\item バナナ
\end{itemize}

\subsection{順序リスト}
\begin{enumerate}
\item 睦月
\item 如月
\item 弥生
\end{enumerate}

\subsection{定義リスト}
\begin{description}
\item[Monday] 月曜日
\item[Tuesday] 火曜日
\item[Wednesday] 水曜日
\end{description}

\subsection{テーブル}

\begin{longtable}[c]{@{}crr@{}}
\toprule\addlinespace
果物 & 値段 & 個数
\\\addlinespace
\midrule\endhead
バナナ & 134円 & 10本
\\\addlinespace
オレンジ & 210円 & 10個
\\\addlinespace
\bottomrule
\end{longtable}

\subsection{画像}
\includegraphics[width=\textwidth]{latex.eps}
